% -----------------------------------------------------------------------------------------------------
\section{$x$ 空间的 DIS 定理（扭度 2 部分）}\label{sec:pdf-dis}

\cref{thm:dis-moments}是关于矩的陈述，\cref{thm:pdf-moments-local}把其中的 $A^{(n)}_i$ 认同为 PDF 的矩；加上矩问题的唯一性（\cref{app:appK}），就得到结构函数扭度 2 部分在 $x$ 空间的因子化；被控制的只是扭度 2 部分，余项只在矩的意义下是 $\order{1/Q^2}$。

\begin{theorem}[DIS 的 $x$ 空间因子化（扭度 2 部分）\tagR]\label{thm:pdf-dis-xspace}
（$R_{\mathrm{OPE}}$：以 OPE 为前提而严格。QCD 中含夸克的微扰 OPE 是全阶论证（\cref{thm:dis-ope}），严格证明只对标量理论与纯 Yang--Mills 的规范不变算符给出；对强子态还须非微扰 OPE（\cref{asm:intro-np-ope}）。）
定义 $F_2^{\mathrm{tw2}}(x,Q^2)$ 为偶数矩满足 $\int_0^1\dd x\,x^{n-2}F_2^{\mathrm{tw2}}=\sum_iC^{(n)}_{2,i}\big(\frac{Q^2}{\mu^2},\as(\mu)\big)A^{(n)}_i(\mu)$（$n=2,4,\dots$）的分布；它由这些矩唯一确定（对强子 PDF 所属的可积函数类由\cref{cor:appK-pdf}，对部分子靶所属的分布类见证明）。则
\begin{equation}\label{eq:pdf-xspace}
F_2^{\mathrm{tw2}}(x,Q^2)/x=\textstyle\sum_i\big(C_{2,i}\otimes f_i\big)(x),
\end{equation}
$C_{2,i}$ 是 Mellin 矩为 $C^{(n)}_{2,i}$ 的唯一分布，构造地等于 $\MSbar$ 减除了共线极点的部分子结构函数。在 $\order{\as}$，记 $L=\ln(Q^2/\mu^2)$，
\begin{equation}\label{eq:pdf-F2nlo}
\frac{F_2^{\mathrm{tw2}}}{x}=\sum_qe_q^2\Big\{(q+\bar q)+\frac{\as}{2\pi}\Big[(C_2^q+L\,P_{qq})\otimes(q+\bar q)+2\,(C_2^g+L\,P_{qg})\otimes g\Big]\Big\}+\order{\as^2},
\end{equation}
\begin{equation}\label{eq:pdf-C2}
\begin{aligned}
C_2^q&=C_F\Big[2\mathcal D_1-\tfrac32\mathcal D_0-(1+z)\ln(1-z)-\tfrac{1+z^2}{1-z}\ln z+3+2z-\big(\tfrac92+\tfrac{\pi^2}3\big)\delta(1-z)\Big],\\
C_2^g&=T_F\Big[\big(z^2+(1-z)^2\big)\ln\tfrac{1-z}{z}-8z^2+8z-1\Big],
\end{aligned}
\end{equation}
% VF-DIS-7
$\mathcal D_k=[\ln^k(1-z)/(1-z)]_+$；纵向结构函数为 $\frac{F_L^{\mathrm{tw2}}}{x}=\frac{\as}{2\pi}\sum_qe_q^2[C_L^q\otimes(q+\bar q)+2C_L^g\otimes g]+\order{\as^2}$，$C_L^q=2C_Fz$，$C_L^g=4T_Fz(1-z)$。$F_2-F_2^{\mathrm{tw2}}$ 的每个固定偶数矩为 $\order{1/Q^2}$（强子：$n\ge n_0$；在\cref{asm:intro-growth}之下一切偶数 $n\ge2$）。
非微扰地位：\cref{asm:intro-wightman}、\cref{asm:intro-composite}、\cref{asm:intro-np-ope}、\cref{asm:intro-af-coeff}、\cref{asm:intro-growth}（MA0、MA3、MA4）。
\uses{thm:dis-moments,thm:pdf-moments-local,lem:prelim-mellin,cor:appK-pdf,prop:appK-kernels,prop:pdf-partonic-dr,asm:intro-wightman,asm:intro-composite,asm:intro-np-ope,asm:intro-af-coeff,asm:intro-growth}
\end{theorem}

\begin{proof}
(1) 由\cref{thm:dis-moments}，对偶数 $n$，$\int_0^1x^{n-2}F_2=\sum_iC^{(n)}_{2,i}A^{(n)}_i+\order{1/Q^2}$（部分子靶：一切偶数 $n\ge2$；强子：$n\ge n_0$，或在\cref{asm:intro-growth}之下一切偶数 $n$）。(2) 由\cref{thm:pdf-moments-local}，偶数 $n$ 时 $A^{(n)}_q=\int_0^1x^{n-1}(q+\bar q)$，$A^{(n)}_g=\int_0^1x^{n-1}g$。(3) \emph{$C_{2,i}$ 的构造。}把 (1)(2) 用于纯维数正规化中的无质量在壳部分子靶 $j$：除 $Q$ 外无别的标度，幂次修正为零；由\cref{prop:pdf-partonic-dr}，$f_{i/j}=Z_{ij}$，故部分子结构函数满足 $\int_0^1z^{n-2}\hat F_{2,j}=\sum_iC^{(n)}_{2,i}Z_{ij}(n)$。令 $C_{2,\cdot}\equiv(\hat F_{2,\cdot}/z)\otimes Z^{-1}$（味道矩阵），它支撑在 $[0,1]$ 上，由\cref{lem:prelim-mellin}其 Mellin 矩恰为 $C^{(n)}_{2,i}$。(4) \emph{唯一性。}每一阶的 $C_{2,i}$ 属于\cref{prop:appK-kernels}(b) 的分布类（$\delta(1-z)$、$\mathcal D_k=[\ln^k(1-z)/(1-z)]_+$ 与在端点至多有 $z^{-1}$ 乘对数奇异性的函数），由该命题，它由偶数矩唯一确定。对强子，$\sum_iC_{2,i}\otimes f_i$ 是满足 $x^{N_1-1}(\cdot)\in L^1$ 的函数，唯一性即\cref{cor:appK-pdf}。(5) 由\cref{lem:prelim-mellin}，$\int_0^1x^{n-1}(C_{2,i}\otimes f_i)=C^{(n)}_{2,i}\int_0^1x^{n-1}f_i$，故\cref{eq:pdf-xspace}右边的偶数矩为 $\sum_iC^{(n)}_{2,i}A^{(n)}_i$，由 (4) 它就是 $F_2^{\mathrm{tw2}}/x$。(6) \emph{$\order{\as}$。}部分子结构函数为 $\hat F_2^q/(ze_q^2)=\delta(1-z)+\frac{\as}{2\pi}(\frac{\mu^2}{Q^2})^\varepsilon[-\frac1\varepsilon P_{qq}+C_2^q+\order\varepsilon]$、$\hat F_2^g/(ze_q^2)=\frac{\as}{2\pi}(\frac{\mu^2}{Q^2})^\varepsilon[-\frac2\varepsilon P_{qg}+2C_2^g+\order\varepsilon]$（胶子道的因子 $2$ 来自每味的夸克与反夸克）。按 (3) 乘以 $Z^{-1}=\delta+\frac{\as}{2\pi}\frac1\varepsilon P+\order{\as^2}$：$(1-\varepsilon L)(-\frac1\varepsilon P+C)+\frac1\varepsilon P=C+LP+\order\varepsilon$，即\cref{eq:pdf-F2nlo}。$F_L$ 由 $F_1$ 的矩同样得到，$C_L=C_2-C_1$ 在这一阶没有极点，与 $\mu$ 无关。(7) $F_2-F_2^{\mathrm{tw2}}$ 的偶数矩为 $\order{1/Q^2}$ 直接来自 (1)。
\end{proof}

两个检验：$C_2^q$、$C_2^g$、$C_L$ 的 Mellin 矩重现\cref{ch:dis}单圈 Wilson 系数的闭式（如 $\int_0^1C_2^q=0$，即 Adler 求和规则；$C^q_{2,2}=\frac16C_F$，$C^g_{2,2}=-\frac14T_F$）；由\cref{thm:pdf-dglap}可验证 $\partial F_2^{\mathrm{tw2}}/\partial\ln\mu^2=\order{\as^2}$（与 $\partial L/\partial\ln\mu^2=-1$ 的项抵消）。定理控制的只是扭度 2 部分与余项的\emph{矩}。余项的逐点大小不受控制：$\sin(Q^2x)$ 的每个矩乘以 $Q^2$ 都有界而 $\sup|\sin(Q^2x)|=1$（\cref{rem:dis-nonuniform}）；OPE 所能给出的最好结果是弱收敛（\cref{cor:dis-weak-xspace}），而“存在矩为给定扭度 2 矩的测度”本身是 Hausdorff 矩问题（\cref{rem:appK-counterexamples}）。在固定 $0<x<1$ 处带 $\order{\Lambda^2/Q^2}$ 余项的逐点 $x$ 空间因子化是\cref{ch:disall}的全阶论证（\cref{thm:disall-factorization}，\tagA），不是本定理的推论。

% -----------------------------------------------------------------------------------------------------
\section{求和规则、方案与正定性}\label{sec:pdf-sumrules}

\begin{proposition}[求和规则 \tagR]\label{prop:pdf-sum-rules}
（$R_{\mathrm{PT}}$：在重整化微扰论中对红外正规化的部分子态严格。）
(a) 数求和规则：$\int_0^1\dd x\,[q(x;\mu)-\bar q(x;\mu)]=N_q$ 对一切 $\mu$ 成立，$N_q$ 是靶的味道 $q$ 的净夸克数（守恒荷的本征值；对部分子靶 $j$，$N_q=\delta_{qj}-\delta_{\bar qj}$）；等价地 $\gamma_{qq}(1)=0$ 在各阶成立。
(b) 动量求和规则：$\sum_i\int_0^1\dd x\,x\,f_i(x;\mu)=1$ 在 $\MSbar$ 中对一切 $\mu$ 精确成立（$T^{++}$ 有限；$N=2$ 反常维度矩阵有本征值 $0$）。
非微扰地位：对强子需\cref{asm:intro-wightman}、\cref{asm:intro-composite}（MA0）。
\uses{thm:pdf-moments-local,prop:pdf-invariance,prop:pdf-lo-kernels,thm:pdf-dglap,asm:intro-wightman,asm:intro-composite}
\end{proposition}

\begin{proof}
(a) \cref{thm:pdf-moments-local}取 $n=1$：$\int_0^1(q-\bar q)=\int_{-1}^1f_q\,\dd\xi=\langle P|[\bar\psi_q\gamma^+\psi_q]_\mu|P\rangle/(2P^+)$。$j^\mu_q=\bar\psi_q\gamma^\mu\psi_q$ 是味道 $q$ 的数守恒流，其荷是好量子数、与标度无关，故该算符不需重整化（单圈即 $\gamma_{qq}(1)=0$，\cref{prop:pdf-lo-kernels}）。由平移不变性，$\int\dd x^-\dd^2\bm x_{\mathrm T}\langle P'|j^+_q(x)|P\rangle=(2\pi)^3\delta^{(3)}(P'-P)\langle P|j^+_q(0)|P\rangle$，而左边等于荷的本征值乘以 $\langle P'|P\rangle$，故 $\langle P|j^+_q(0)|P\rangle=2P^+N_q$。
(b) \cref{thm:pdf-moments-local}取 $n=2$ 并对各味求和：$\sum_q\int_0^1x(q+\bar q)+\int_0^1xg=[\sum_q\langle O^{(2)}_q\rangle+\langle O^{(2)}_g\rangle]_\mu/(2(P^+)^2)$。规范不变对称能动张量 $T^{\mu\nu}=\sum_q\frac{\ii}4\bar\psi_q(\gamma^\mu\overleftrightarrow D^\nu+\gamma^\nu\overleftrightarrow D^\mu)\psi_q-G^{\mu\alpha}_aG^\nu{}_{\alpha a}+\frac14g^{\mu\nu}G^2$（规范固定项与鬼场项是 BRST 恰当的，在物理态之间不贡献）的 $++$ 分量在前向矩阵元中等于 $\sum_qO^{(2)}_q+O^{(2)}_g$（夸克部分把导数移到右边，胶子部分见\cref{prop:pdf-invariance}(d)）。能动张量守恒，其荷为四动量，故有限、不需重整化（\cite{Collins1984}），且同 (a) 一样 $\langle P|T^{\mu\nu}|P\rangle=2P^\mu P^\nu$，于是左边等于 $\langle P|T^{++}|P\rangle/(2(P^+)^2)=1$。用 DGLAP 的语言，这等价于 $N=2$ 的全阶反常维度矩阵各列之和为零；单圈时即\cref{prop:pdf-lo-kernels}中的动量求和规则（\cref{thm:pdf-dglap}）。
\end{proof}

对强子，$N_q$ 是价夸克数（质子 $N_u=2$，$N_d=1$）；强子\emph{结构函数}的动量求和规则另需增长条件\cref{asm:intro-growth}，它与低阶矩的色散表示等价（\cref{prop:dis-low-moments}）。

\begin{remark}[解析延拓与 $P^\pm_{\mathrm{NS}}$]\label{rem:pdf-ns-plusminus}
定域算符只给出整数 $N$ 处的反常维度，电磁 DIS 中只出现偶数 $N$；$x$ 空间的核是某个解析延拓的 Mellin 逆变换。由偶数 $N$ 出发的延拓（C 偶，$q+\bar q$）与由奇数 $N$ 出发的延拓（C 奇，$q-\bar q$）一般是\emph{不同}的函数；在多项式有界的条件下各自唯一（\cref{prop:appK-continuation}），例如 $S_1(N)\mapsto\psi(N+1)+\gamma_E$，而加上 $c\sin(\pi N/2)$ 因指数增长被排除。领头阶 $qq$ 道两种延拓相同；从次领头阶起 $P^+_{\mathrm{NS}}-P^-_{\mathrm{NS}}=2P^V_{q\bar q}\ne0$，色因子为 $C_F(C_F-\frac{C_A}2)$\cite{CurciFurmanskiPetronzio1980}。以 $(\as/2\pi)^2$ 为展开参数，$\int_0^1P^{V(1)}_{q\bar q}\,\dd x=0.71931\,C_F(C_F-\frac{C_A}2)\ne0$，而 $\int_0^1P^-_{\mathrm{NS}}=0$（价夸克数守恒），所以 $P^+_{\mathrm{NS}}$ 的第一矩不为零。“DGLAP 方程与 OPE 重整化群方程在所有阶以单一的核等价”是错的：两者等价，但非单态有两个核，分别对应两种 C 宇称。
% VF-REF-5
\end{remark}

\begin{remark}[方案变换]\label{rem:pdf-schemes}
$\MSbar$ 只是约定。对有限变换 $f'=Z\otimes f$（$Z=\mathbb 1+\frac{\as}{2\pi}Z_1+\cdots$，与靶无关），系数函数变为 $C'=C\otimes Z^{-1}$，可观测量不变，演化核变为 $\mathcal P'=Z\mathcal PZ^{-1}+\frac{\dd Z}{\dd\ln\mu^2}Z^{-1}$，即 $P'^{(1)}=P^{(1)}+Z_1\otimes P^{(0)}-P^{(0)}\otimes Z_1-\frac{\beta_0}2Z_1$（$\mathcal P=\frac{\as}{2\pi}P^{(0)}+(\frac{\as}{2\pi})^2P^{(1)}$）。DIS 方案选 $Z_1$ 使 $F_2$ 保持部分子模型的形式：$q^{\mathrm{DIS}}_i=q_i+\frac{\as}{2\pi}[C_2^q\otimes q_i+C_2^g\otimes g]$，并按动量守恒令 $g^{\mathrm{DIS}}=g-\frac{\as}{2\pi}[C_2^q\otimes\Sigma+2n_fC_2^g\otimes g]$（参见\cite{DiemozEtAl1988}），于是 $P^{(1),\mathrm{DIS}}_{\mathrm{NS}}=P^{(1),\MSbar}_{\mathrm{NS}}-\frac{\beta_0}2C_2^q$，$C_L$ 在 $\order{\as}$ 与方案无关。胶子极化的平均约定（\cref{def:pdf-gluon}）也是方案的一部分。DIS 方案是合法的方案选择；但若把 PDF \emph{定义}为“$F_2$ 中的那个量”，其他过程中的普适性就成了定义而非定理。
% VF-DIS-11
\end{remark}

\begin{conjecture}[$\MSbar$ PDF 的正定性 \tagC]\label{conj:pdf-positivity}
存在一个（非微扰的）标度 $\mu_0$，使得对一切 $\mu\ge\mu_0$，强子的 $\MSbar$ 夸克、反夸克与胶子分布非负：$q(x;\mu),\,\bar q(x;\mu),\,g(x;\mu)\ge0$（$0<x<1$）。
\uses{def:pdf-quark,def:pdf-gluon}
\end{conjecture}

\paragraph{地位与缺口}
没有一般的证明，而更强的命题是错的。\emph{支持：}Candido、Forte 与 Hekhorn\cite{CandidoForteHekhorn2020}论证次领头阶 $\MSbar$ PDF 在微扰区域为正；Candido、Forte、Giani 与 Hekhorn\cite{Candido:2023ujx}同意 PDF 在足够低的非微扰标度会变负，并保守地半定量估计 $Q^2\gtrsim5\,\mathrm{GeV}^2$ 以上正定性得以继承。\emph{反对一般定理：}Collins、Rogers 与 Sato\cite[第~VI~节]{CollinsRogersSato2022}指出上述论证在“裸 PDF 为正”这一步失效——用维数正规化同时处理紫外（$n<4$）与共线（$n>4$）发散时，在 $n<4$ 为正的紫外贡献延拓到 $n>4$ 后可以变负。\emph{“在所有标度上非负”是错的：}Yukawa 模型中单圈 $\MSbar$ 分布为 $f_{q/p}(\xi;\mu)=a_\lambda(1-\xi)[\ln\frac{\mu^2}{\Delta(\xi)}-1+\frac{(m_q+\xi m_p)^2}{\Delta(\xi)}]$，$\Delta(\xi)=\xi m_s^2+(1-\xi)m_q^2-\xi(1-\xi)m_p^2$，$a_\lambda=\frac{\lambda^2}{16\pi^2}$\cite[式~(55)]{CollinsRogersSato2022}；
% VF-FRO-10
对 $m_q,m_p,m_s=0.3,1.0,1.5\,\mathrm{GeV}$，$\mu=0.8\,\mathrm{GeV}$ 时恰在 $\xi\in(0.264,1)$ 上 $f<0$，$\mu=1.2\,\mathrm{GeV}$ 时在 $\xi>0.567$ 上 $f<0$，而 $f$ 对一切 $\xi$ 为正当且仅当 $\mu\ge\mu_*=m_s\exp\{\frac12[1-(m_q+m_p)^2/m_s^2]\}=1.699\,\mathrm{GeV}$（这些数值是我们对该式的直接计算）。QCD 中也有初等的例子：单圈匹配给出 $\MSbar$ 重夸克分布 $f_h(x;\mu)=\frac{\as}{2\pi}\ln\frac{\mu^2}{m_h^2}(P_{qg}\otimes g)(x)+\order{\as^2}$（\cref{ex:pdf-quark-in-gluon}），$\mu<m_h$ 且 $g>0$ 时为负。\emph{循环论证的警示：}领头阶 DGLAP 演化保持正定性，并不证明初始条件为正，而高阶的正定性保持依赖于方案。因此猜想只针对 $\mu\ge\mu_0$，而 $\mu_0$ 是非微扰的；本书只陈述这些事实，不作进一步的裁决，详细讨论（含 Yukawa 曲线）见\cref{ch:front}。强子 $\MSbar$ PDF 本身的存在依赖于\cref{asm:intro-wightman}与\cref{asm:intro-composite}；此缺口与幂次计数无关，不涉及 E-S 至 E-W 型的误差估计。升级途径：一个证明或一个 $\mu\ge\mu_0$ 区域的反例，以及 $\mu_0$ 的非微扰确定。
